% =====================================================
% 题型：立体几何单选题 (q_solid_cuboid_skew_angle.tex)
% =====================================================
% =====================================================
% 难度：★★ (中档)
% =====================================================
\newcommand{\qSolidCuboidSkewAngleStem}{%
  在长方体 $ABCD - A_1B_1C_1D_1$ 中，$AB = AD = 2\sqrt{3}$，$AA_1 = 2$，则直线 $AD_1$ 与 $BD$ 所成角的余弦值为%
}

\newcommand{\qSolidCuboidSkewAngleOptions}{%
  \mcoptions[4]{$\dfrac{\sqrt{3}}{3}$}{$\dfrac{\sqrt{3}}{4}$}{$\dfrac{\sqrt{6}}{3}$}{$\dfrac{\sqrt{6}}{4}$}%
}

\newcommand{\qSolidCuboidSkewAngleFigure}[1][1.0]{%
  \begin{tikzpicture}[scale=#1, >=stealth]
    % Coordinates for Cuboid ABCD-A1B1C1D1
    \coordinate (D) at (0,0);
    \coordinate (C) at (3.0,0);
    \coordinate (A) at (-1.2,-1.2);
    \coordinate (B) at (1.8,-1.2);
    
    \coordinate (D1) at (0,2);
    \coordinate (C1) at (3.0,2);
    \coordinate (A1) at (-1.2,0.8);
    \coordinate (B1) at (1.8,0.8);
    
    % Visible edges
    \draw[thick] (A) -- (B) -- (C) -- (C1) -- (D1) -- (A1) -- cycle;
    \draw[thick] (A1) -- (B1) -- (C1);
    \draw[thick] (B) -- (B1);
    \draw[thick] (A) -- (A1);
    
    % Hidden edges
    \draw[thick, dashed] (A) -- (D) -- (C);
    \draw[thick, dashed] (D) -- (D1);
    
    % Diagonals defined in the question (blue dashed)
    \draw[thick, dashed, blue] (A) -- (D1);
    \draw[thick, dashed, blue] (B) -- (D);
    
    % Points
    \foreach \p in {A,B,C,D,A1,B1,C1,D1} {
      \fill (\p) circle (1.2pt);
    }
    
    % Labels
    \node[below left] at (A) {\small $A$};
    \node[below right] at (B) {\small $B$};
    \node[right] at (C) {\small $C$};
    \node[above left] at (D) {\small $D$};
    \node[left] at (A1) {\small $A_1$};
    \node[right] at (B1) {\small $B_1$};
    \node[right] at (C1) {\small $C_1$};
    \node[above left] at (D1) {\small $D_1$};
    
    \ifteacher
      % Coordinate system
      \draw[->, gray, dashed, thin] (D) -- (4.0,0) node[right] {\small $y$};
      \draw[->, gray, dashed, thin] (D) -- (-2.0,-2.0) node[below left] {\small $x$};
      \draw[->, gray, dashed, thin] (D) -- (0,3.0) node[above] {\small $z$};
      
      % Teacher-only translation helper lines (gray dashed)
      \draw[thick, dashed, gray] (B) -- (C1);
      \draw[thick, dashed, gray] (C1) -- (D);
    \fi
  \end{tikzpicture}%
}

\newcommand{\qSolidCuboidSkewAngleAnswer}{D}

\newcommand{\qSolidCuboidSkewAngleSolution}{%
  在长方体 $ABCD - A_1B_1C_1D_1$ 中，因为 $AB \parallel D_1C_1$ 且 $AB = D_1C_1$，
  
  所以四边形 $ABC_1D_1$ 是平行四边形，故
  \[ AD_1 \parallel BC_1. \]
  
  从而有，直线 $AD_1$ 与 $BD$ 所成的角即为直线 $C_1B$ 与 $BD$ 所成的角 $\angle C_1BD$（或其补角）。
  
  连接 $C_1D$．因为 $AB = AD = 2\sqrt{3}$ 且 $AA_1 = 2$，
  
  在 Rt$\triangle C_1CB$ 中，因为 $CC_1 = AA_1 = 2$，$BC = AD = 2\sqrt{3}$，
  
  所以对角线
  \[ BC_1 = \sqrt{BC^2 + CC_1^2} = \sqrt{(2\sqrt{3})^2 + 2^2} = 4. \]
  
  同理，在 Rt$\triangle C_1CD$ 中，因为 $CD = AB = 2\sqrt{3}$，$CC_1 = 2$，
  
  所以对角线
  \[ C_1D = \sqrt{CD^2 + CC_1^2} = \sqrt{(2\sqrt{3})^2 + 2^2} = 4. \]
  
  在底面 Rt$\triangle BAD$ 中，因为 $AB = AD = 2\sqrt{3}$，
  
  所以对角线
  \[ BD = \sqrt{AB^2 + AD^2} = \sqrt{(2\sqrt{3})^2 + (2\sqrt{3})^2} = 2\sqrt{6}. \]
  
  在等腰 $\triangle C_1BD$ 中，因为 $BC_1 = C_1D = 4$，$BD = 2\sqrt{6}$，
  
  由余弦定理得：
  \[ \cos\angle C_1BD = \frac{BC_1^2 + BD^2 - C_1D^2}{2 \times BC_1 \times BD} = \frac{4^2 + (2\sqrt{6})^2 - 4^2}{2 \times 4 \times 2\sqrt{6}} = \frac{24}{16\sqrt{6}} = \frac{\sqrt{6}}{4}. \]
  
  因为 $\dfrac{\sqrt{6}}{4} > 0$，
  
  所以直线 $AD_1$ 与 $BD$ 所成角的余弦值为 $\dfrac{\sqrt{6}}{4}$。
  
  故选 \textbf{D}。
}

\newcommand{\qSolidCuboidSkewAngleDiagnosis}{%
  本题考查异面直线所成角的两种基本方法：平移法与坐标法。
  学生若只盯着原图中的 $AD_1$ 与 $BD$，容易误以为无法构造平面角；
  关键是发现 $AD_1\parallel BC_1$，从而把异面直线夹角转化为
  $\angle C_1BD$。随后在 $\triangle C_1BD$ 中用余弦定理即可。
}

\newcommand{\qSolidCuboidSkewAngleTeachingNotes}{%
  \item \textbf{几何法核心}：发现 $AD_1\parallel BC_1$，将异面直线夹角转化为 $\angle C_1BD$，再在等腰三角形 $\triangle C_1BD$ 中用余弦定理。
  \item \textbf{坐标法拓展}：本题亦可建立空间直角坐标系。以 $D$ 为原点，分别以 $DA,DC,DD_1$ 为 $x,y,z$ 轴正方向，设 $D(0,0,0)$，则可快速求出两向量方向向量，计算夹角余弦值。
}
